\section{Plan de pruebas y validación}\label{sec:pruebas}

\subsection{Validación de restricciones}\label{sec:validador}
El módulo \texttt{validator.py} verifica cada solución sin usar las estructuras internas de los algoritmos. Recalcula: (1) que cada vuelta aparezca una sola vez, asignada o no atendida; (2) que el bus sea de la ruta; (3) la ventana; (4) la jornada; (5) que la duración coincida con $p_j(s_j)$; (6) que el almuerzo sea único, de 120\,min y dentro de su ventana; y (7) que no haya solapamientos. Cada ejecución desde la línea de comandos lo aplica y termina con error si encuentra violaciones.

\subsection{Casos de prueba}
{\small
\begin{longtable}{L{2.9cm}L{6.4cm}L{5.4cm}}
\caption{Plan de pruebas de la Entrega 1 (automatizado con pytest).}\label{tab:pruebas}\\
\toprule
Tipo & Caso & Resultado esperado \\
\midrule\endfirsthead
\toprule
Tipo & Caso & Resultado esperado \\
\midrule\endhead
\bottomrule\endfoot
Básico & Instancia manual con LS y LPT en T0, T1 y T2 & Soluciones válidas que atienden las 8 vueltas \\
Básico & Cargas máximas de la instancia manual & T0: 270 (LS y LPT); T2: 329.0 (LS) y 333.2 (LPT) \\
Básico & Asignación de cada vuelta en el ejemplo manual & Coincide con las trazas de la sección~\ref{sec:manual} \\
Límite & Ventana de ancho cero ($r_j = d_j$) & Salida exactamente en $r_j$ \\
Límite & Vuelta que termina exactamente a las 22:00 & Se acepta \\
Límite & Un solo bus ($m = 1$) & Solución válida \\
Límite & Duraciones iguales sin tráfico & LPT coincide con LS \\
Adverso & 3 vueltas simultáneas y 1 bus & 2 vueltas no atendidas, sin violaciones \\
Adverso & Peor caso clásico de LPT $\{3, 3, 2, 2, 2\}$\,h con $m = 2$ & LPT obtiene 7\,h (el óptimo es 6\,h) \\
Ventanas & 80 vueltas con ventanas y tráfico T2 & Toda salida dentro de su ventana \\
Tráfico & Duración asignada en la instancia manual & Igual a $p_j(s_j)$ del perfil \\
Tráfico & Propiedad FIFO en pares aleatorios & $s_1 \le s_2 \Rightarrow$ llegada$_1 \le$ llegada$_2$ \\
Tráfico & Cálculo de $p_{J1}(06{:}00)$ en T2 y media de T1 & 103.67\,min; media de T1 = 1 \\
Tráfico & Carga total con T2 frente a T0 & Mayor con T2 \\
Almuerzo & Vuelta que choca con un almuerzo fijo & Queda no atendida \\
Almuerzo & Vuelta que choca con un almuerzo flexible & El almuerzo se desplaza y la vuelta se atiende \\
Almuerzo & Almuerzo de cada bus & Dentro de su ventana y de 120\,min \\
Validación de restricciones & Solución alterada a propósito & El validador la rechaza \\
Validación de restricciones & Cota inferior frente a $\Lmax$ & $LB \le \Lmax$ \\
Validación de restricciones & 400 vueltas y 60 buses & Solución válida; todas las vueltas contabilizadas \\
Validación de datos & Lectura del Excel del proyecto & 35 rutas, 4\,733 vueltas, coinciden con la regla general \\
\end{longtable}}

La suite completa (\texttt{tests/}) contiene 36 pruebas. En el equipo del grupo se aprobaron 35 y se omitió 1: la prueba del prototipo CP-SAT de \texttt{test\_referencia.py}, que solo se ejecuta si está instalado el paquete opcional \texttt{ortools}. Esa prueba corresponde a los mecanismos de referencia previstos para una etapa posterior y no afecta a LS ni a LPT.

\captura{pytest.png}{Ejecución de las pruebas automatizadas (\texttt{python -m pytest -q}): 35 pruebas aprobadas y 1 omitida.}{fig:cap-pytest}

\subsection{Verificación de las salidas}
\begin{enumerate}
  \item \textbf{Factibilidad:} el validador recalcula todas las restricciones en cada ejecución.
  \item \textbf{Cálculo a mano:} la instancia de la sección~\ref{sec:manual} y la calculadora del Excel coinciden con el programa.
  \item \textbf{Calidad:} la carga máxima se compara con la cota inferior $LB$ de las vueltas atendidas.
  \item \textbf{Instancias nuevas:} cualquier instancia escrita en el formato de la sección~\ref{sec:sistema} se resuelve y valida con \texttt{python -m bus\_sched run --instance archivo.json}.
\end{enumerate}
La comparación con la solución exacta en instancias pequeñas corresponde a los mecanismos de referencia de la sección~\ref{sec:referencia}, previstos para una etapa posterior.
